\section{Methodology}
\label{sec:method}

We describe a three-stage pipeline that transforms web page text into a two-dimensional vibe map and detail the evaluation protocol we use to assess whether vibe space captures information distinct from content space.

\subsection{Data}
\label{sec:data}

We draw from the \cfour (Colossal Clean Crawled Corpus) small sample~\cite{raffel2020exploring}, which contains 10{,}000 pre-extracted web page texts hosted on HuggingFace.
After filtering to pages with at least 200 characters (removing 447 pages), we randomly sample 500 pages with seed 42.
Three pages fail during LLM processing due to transient API errors, yielding a final working sample of $N{=}497$.

\subsection{Vibe Generation}
\label{sec:vibe_gen}

For each page, we truncate the text to 2{,}000 characters and prompt \gptfour with the following:

\begin{quote}
\small
\textbf{System:} ``You are a cultural commentator who describes the vibe, aesthetic, and energy of things using contemporary internet language. Be specific, evocative, and concise (1--3 sentences).''

\textbf{User:} ``Here is the content of a web page: [TEXT]. What is this web page giving? Start your answer with `it's giving'.''
\end{quote}

We use temperature 0.7 and a maximum of 150 tokens.
The prompt is designed to elicit descriptions that foreground aesthetic character---era, tone, community type, cultural energy---rather than topical content.
The ``it's giving'' framing draws on Gen~Z internet vernacular, encouraging the model to adopt a culturally fluent register.

\subsection{Embedding and Projection}
\label{sec:embedding}

We embed both the raw page texts and the vibe descriptions using \mpnet~\cite{reimers2019sentence}, a 768-dimensional sentence transformer.
We then project both sets of embeddings to two dimensions using UMAP~\cite{mcinnes2018umap} with $n_{\text{neighbors}}{=}15$, $\text{min\_dist}{=}0.1$, and cosine distance---standard settings that balance local and global structure preservation.

\subsection{Evaluation Protocol}
\label{sec:eval}

We evaluate four hypotheses about the relationship between vibe space and content space.

\para{H1: Vibe diversity.}
We measure whether the LLM produces meaningfully varied descriptions by computing the number of unique descriptions, vocabulary richness (type-token ratio), and mean description length.

\para{H2: Vibe $\neq$ topic.}
We cluster both embedding spaces with KMeans~\cite{arthur2007kmeans} for $k \in \{5, 10, 15, 20\}$ and compute the Adjusted Rand Index (\ari)~\cite{hubert1985ari} between the two clusterings.
Low \ari indicates that vibe space groups pages differently from content space.
We assess significance with a permutation test (10{,}000 permutations).

\para{H3: Cluster quality.}
We compute silhouette scores~\cite{rousseeuw1987silhouettes} for each space and compare them using bootstrap 95\% confidence intervals (1{,}000 resamples).

\para{H4: Cross-domain discovery.}
For each page, we identify its $K{=}5$ nearest neighbors in both vibe space and content space (using high-dimensional cosine similarity).
We measure (a)~the overlap between the two neighbor sets, and (b)~the raw-text cosine similarity of vibe-space neighbors.
If vibe space enables cross-domain discovery, vibe-space neighbors should have substantially lower content similarity than content-space neighbors.
We report the Mann--Whitney $U$ statistic and Cohen's $d$ effect size.
